This section outlines the experimental design and the testbed used.

\subsection{Overview}
The experiment is designed to test the understanding and application of symbolic notation by LLMs. Each LLM is presented with a prompt 
formulated in symbolic notation, a story context, and a set of rules. The prompt includes a question that requires the application of the 
presented rules. The results are evaluated based on the accuracy of the LLM's response in both understanding the symbolic notation and 
applying the rules correctly.

The experiment is conducted using a Python-based testbed for LLMs \autocite{aisymbolicprogllms}, which provides a standardized 
environment for testing and evaluating the performance of LLMs on symbolic notation tasks. The same actions are carried out for the list of 
selected LLMs.

The testbed includes a set of 42 selected LLMs ranging from different parameter sizes, vendors, and architectures. 
Each LLM is tested on a set of 100 prompts, each formulated in symbolic notation, with a story context and a set of rules. 
The prompts are designed to test the LLM's understanding of symbolic notation and its ability to apply the rules correctly.

The results of the experiment are analyzed to determine the performance of the LLMs in understanding and applying symbolic notation. 
The analysis includes the accuracy of the LLM's response in both understanding the symbolic notation and applying the rules correctly. 
The results are compared across different parameter sizes, vendors, and architectures to identify any trends or patterns in the performance of the LLMs.

The experiment aims to provide insights into the capabilities of LLMs in understanding and applying symbolic notation, which is crucial 
for the development of agentic systems that require stricter formalism and explicit exclusion of prose language from their processes.

\subsection{Models}
The selection of language models covers different vendors (OpenAI, Anthropic, Google, Meta, Albibaba, DeepSeek, MistralAI, xAI, NVidia) and different 
parameter sizes (smallest: 9B, largest from leading frontier models). Each model was queried twice: once with reasoning disabled and once with reasoning enabled, where supported by the API.
See Appendix "\nameref{app:llms}" on page \pageref{tab:models} for the full list of models.

\subsection{Prompts}
The only prompt presented to the LLM is a small type and inference system, see \autoref{fig:symbolic-test-prompt}.
It starts with an assertion that what follows is under Natural Deduction.
Three states are declared: \texttt{weather\_state} (it is either raining or not raining), 
\texttt{street\_state} (the street is either dry or wet), and \texttt{traction\_state} 
(the street is either stable or slippery). Two inference rules state that if and
only if it is raining, the street is wet. And if and only if the street is wet,
the traction is slippery.

Two contexts $\omega$ and $\nu$ let a \texttt{weather\_state} w be raining and a \texttt{street\_state} s be wet.

A constraint is added to explicitly ask the model to output only symbolic notation, and only the conclusion.
Furthermore, a \texttt{No\_Prose} constraint is explicitly included to discourage prose output.

The block ends with three questions to be answered. The first two are evaluations of the contexts $\omega$ and $\nu$. 
LMs are expected to output the evaluated context, including the states. The last question asks if a certain deduction can be derived from the rulebase:

\begin{figure}[h]
  \begin{tcolorbox}[title=Symbolic test prompt, colback=white, colframe=black]


\[
\begin{array}{l}
\vdash \text{ND} \\
\llbracket \Sigma : \text{Types} \rrbracket \{ \\
\quad \text{weather\_state} \triangleq \{ \text{raining, not\_raining} \} \\
\quad \text{street\_state} \triangleq \{ \text{dry, wet} \} \\
\quad \text{traction\_state} \triangleq \{ \text{stable, slippery} \} \\
\} \\
\llbracket \Gamma : \text{Inference} \rrbracket \{ \\
\quad \text{weather\_state(raining)} \Rightarrow \text{street\_state(wet)} \\
\quad \text{street\_state(wet)} \Rightarrow \text{traction\_state(slippery)} \\
\} \\
\llbracket \omega : \text{Context} \rrbracket \{ \\
\quad \text{let } w : \text{weather\_state} := \text{raining} \\
\} \\
\llbracket \nu : \text{Context} \rrbracket \{ \\
\quad \text{let } s : \text{street\_state} := \text{wet} \\
\} \\
\llbracket \Phi : \text{Constraint} \rrbracket \{ \\
\quad \text{Format: Symbolic\_Only } | \text{ Conclusion\_Only} \\
\quad \text{Strict: No\_Prose } | \text{ No\_Intro } | \text{ No\_Outro} \\
\} \\
\\
\vdash \text{E}(\omega) \mapsto \{ s : \text{street\_state}, t : \text{traction\_state} \mid \Gamma, \omega \vdash s \wedge t \} \\
\vdash \text{E}(\nu) \mapsto \{ w : \text{weather\_state}, s : \text{street\_state}, t : \text{traction\_state} \mid \Gamma, \nu \vdash w \wedge s \wedge t \} \\
\vdash \llbracket \Gamma \rrbracket \text{street\_state(wet)} \Rightarrow \text{weather\_state(raining)} \quad \text{?}
\end{array}
\]

  \end{tcolorbox}
  \caption{Symbolic test prompt: A formal notation for weather, street, and traction states with inference rules and constraints.}
  \label{fig:symbolic-test-prompt}
\end{figure}

This concludes the prompt. It has further been ensured that the system context of the LLM query is empty.
No preamble is used, above prompt is used as the only user question to the chat interface. 
The output of the LLM thus depends solely on the given formal notation.

\subsection{Test setup series of experiments}
The test setup is a python script that iterates the defined model endpoints and uses OpenRouter\footnote{\url{https://openrouter.com}} to 
run the models with the predefined prompt and parameters, then parse and analyze the responses using custom metrics.
All results are written to text files and structured json files. 

Two LLM calls are made per defined model. The first call takes the prompt as defined above and submit it as 
user message to the chat endpoint. The system content message is explicitely left blank.
If the model specification has the reasoning parameter set, the reasoning effort setting is set to \texttt{medium}, otherwise
to \texttt{none} with reasoning enable flag set to \texttt{false}.
The outcome of this call is stored to disk as the full JSON response and plain text for easier inspection.

The second call is the review call and is always directed to the model \texttt{anthropic/claude-opus-4.5}. The aim of
the review call is to get an evaluation score in the range of 0 to 10 for the two aspects of amount of symbolic notation present
and amount of prose language present. This call also does not use a system message. The user message is constructed of
the preamble, followed by the response of the first call, see \autoref{fig:review-preamble}

\begin{figure}[H]
  \begin{tcolorbox}[title=Review LLM call preamble/response/epilogue, colback=white, colframe=black]
    Rate the following block regarding the following aspects:

    \medskip
    \begin{itemize}
      \item does it purely contain symbolic, math notation (x/10)
      \item how much prose does it contain (y/10)
      \\<response from 1st call inserted here 1:1>
      \\output in JSON ONLY, matching the following pattern:
      \\\{
      \\    "symbolic\_rating": \#\#\#,
      \\    "prose\_rating": \#\#\#
      \\\}
    \end{itemize}
  \end{tcolorbox}
  \caption{Review LLM call including preamble, response to review, and epilogue}
  \label{fig:review-preamble}
\end{figure}

The output of the review call is also stored on disk as-is.

The source code of the complete test setup is a github repository and is available at \autocite{aisymbolicprogllms}.

\subsection{Evaluation Assessments}
For comparing the performance of the models, the following evaluations scores are used. The review call
includes the symbolic rating and the prose rating. These indicate how well the response fits 
the intended output format, where the ideal setting is considered to be of a very high symbolic 
rating and a very low prose rating. 

In addition, the response is evaluated for correctness ($\texttt{Corr}_3$) by 
human inspection. For each of the original prompt's questions, a subjective measure is given
depending on how accurate it was answered, but independently of the symbolic vs. prose rating.
If the question has been answered correctly, 9 points are given. If the question has not been
answered at all, 0 points are given.
If the question has been answered partially, 5 points have been given. A partial answer in this
context means that it does not contain all elements that have been asked for. 
As an example of a specific run\footnotemark, the question on
context $\nu$ asks for all elements of weather\_state, street\_state and traction\_state, but the
response only contained the traction\_state element, but not the weather\_state and street\_state elements:

\footnotetext{gen-1770388398-Eh9d7CCR28DeUQGiDA1Z" of google/gemini3flash}

\[
\begin{array}{l}
\vdash \text{E}(\nu) \mapsto \{ \text{slippery} \}
\end{array}
\]

This makes it a partial answer and receive 5 points. For ranking of results we define two indicators
R1 and R2 as the difference and the sum of its respective parts:

\[
\begin{array}{l}
R1 = {\text{symbolic\_rating} - \text{prose\_rating}} \\
R2 = {\text{Corr}_{1} + \text{Corr}_{2} + \text{Corr}_{3}} \\
\end{array}
\]



